\documentclass[10pt,a4paper]{amsart}
\usepackage[margin=19mm]{geometry}
\usepackage{amsmath,amssymb,mathtools}
\usepackage[scr=rsfs,cal=euler]{mathalfa}
\usepackage{xfrac}
\usepackage[unicode,hidelinks,bookmarks=false]{hyperref}
\newcommand{\ie}{i.e.\ }
\newcommand{\cf}{cf.\ }
\newcommand{\resp}{respectively\ }
\newcommand{\Subsection}{\subsection}
\providecommand{\nobreakdash}{\nobreak\hbox{-}}
\setlength{\emergencystretch}{2em}
\multlinegap0pt
\usepackage{bm}
\usepackage{bbm}
\usepackage{dsfont}
\usepackage{xspace}
\usepackage{cancel}
\usepackage{mathtools}
\usepackage{amsthm}
\makeatletter
\@ifundefined{theorem}{\newtheorem{theorem}{Theorem}}{}
\makeatother
\def\pmod #1{\ ({\rm{mod}}\ #1)}
\title[Note on a problem of S\'ark\"ozy on multiplicative represent]{Note on a problem of S\'ark\"ozy on multiplicative representation functions}
\author{Yuchen Ding}
\date{Source: arXiv:2504.09545v2 (CC BY 4.0)}
\begin{document}
\maketitle

\noindent\emph{Source and licence.} Note on a problem of S\'ark\"ozy on multiplicative representation functions. Version arXiv:2504.09545v2 (CC BY 4.0), distributed under the Creative Commons Attribution 4.0 International licence, as recorded in the arXiv licence metadata for this exact version and shown on the abstract page https://arxiv.org/abs/2504.09545v2. Full rights notice: licenses/arxiv-2504.09545-CC-BY-4.0.txt.\par\smallskip

\noindent\textit{Editorial scope.} Counted unit: the Theorem whose source label is \texttt{thm2} in the article (source0.tex lines 141--156, source label \texttt{thm2}), delivered together with the article's own complete original proof of it (source0.tex lines 159--313, one \texttt{proof} environment of 155 lines). No mathematics is written, repaired or re-derived: statement and proof are the article's own text line for line. Required context retained uncounted: none. Every definition, display and label of those lines is retained unchanged. Omitted from this package and not redistributed: 1--140, 157--158, 314--347. Nothing inside a retained excerpt is omitted or altered: every retained body is delivered exactly as the article's lines are written, so no omission receipt of any kind accompanies this package. Citations: the retained excerpts carry no citation command at all, so no bibliography entry of the article is transcribed and no reference is redistributed. Reference adaptation: none was needed and none was made; every reference inside the delivered unit resolves inside it, so no unresolved reference or citation is produced. Printed slips kept literal: no printed word, symbol, hypothesis or number of the counted statement or of its proof was altered. Layout and class: the article's own class is \texttt{amsart} with packages including amsmath, amsfonts, amssymb, amsthm, latexsym, enumerate, url, cases, mathrsfs, hyperref, etoolbox; the wrapper is the lane's pinned \texttt{amsart} editorial wrapper at 10pt on A4, which restates the theorem-like environments the retained text needs and adds the article's own macro definitions that the retained text needs (\texttt{\textbackslash pmod}), with extra wrapper packages (bm, bbm, dsfont, xspace, cancel, mathtools). The delivered numbering is the unit's own. Assets: no retained span contains a figure environment, a graphics-inclusion command, a PDF-page inclusion, a table or a TikZ picture; no image file is copied into this package, the package declares no asset dependency and no image file is redistributed with it. Licence: version arXiv:2504.09545v2 of this article is distributed under the Creative Commons Attribution 4.0 International licence, as recorded in the arXiv licence metadata for this exact version and shown on the abstract page https://arxiv.org/abs/2504.09545v2. A full rights notice is delivered at licenses/arxiv-2504.09545-CC-BY-4.0.txt. Characters: every retained excerpt is pure ASCII, so the delivered engine, XeLaTeX with its default Latin Modern text font, reports no missing character.
\begin{theorem}\label{thm2}
Let $b,c,e$ and $f$ be integers. Then we have
\begin{align}\label{eq4}
\limsup_{n\rightarrow\infty}|d(\mathcal{A},bn+c)-d(\mathcal{A},en+f)|=\infty
\end{align}
for any infinite $\mathcal{A}\subset \mathbb{N}$ 
if and only if one of the following cases holds:

I. $b=c=e=0$ and $f\neq 0$; or $b=c=f=0$ and $e\neq 0$;

II. $b=0$, $e\neq 0$, and $c=0$; or $b\neq 0$, $e= 0$, and $f=0$;

III. $b=0$, $e\neq 0$, $c\neq 0$, and $e|f$; or $b\neq 0$, $e= 0$, $f\neq 0$, and $b|c$;

IV. $be \neq 0$, $bf\neq ec$, and $b|c$; or $be \neq 0$, $bf\neq ec$, and $e|f$.
\end{theorem}

\begin{proof}[Proof of Theorem \ref{thm2}]
{\bf Necessities.}
To prove the necessities, we separate them into a few cases.

{\bf Case 1.} $b=e=0$.

It is clear that the situation $b=c=e=f=0$ does not satisfy our requirement. Moreover, at least one of $c$ and $f$ equals zero. Otherwise, we have
\begin{align*}
\limsup_{n\rightarrow\infty}|d(\mathcal{A},c)-d(\mathcal{A},f)|\le |c|+|f|<\infty
\end{align*}
for any infinite $\mathcal{A}\subset \mathbb{N}$ which is a contradiction. ({\bf Case 1.} corresponds with I.)

{\bf Case 2.} $b=0$ and $e\neq 0$; or $b\neq 0$ and $e= 0$.

We only consider the case $b=0$ and $e\neq 0$ since the other situation will be treated similarly. In this case, $c=0$ is admissible since $d(\mathcal{A},en+f)<\infty$ for all sufficiently large $n$ as $en+f\neq 0$ for $n>-f/e$. We are left over to consider the situation $b=0,e\neq 0$, and $c\neq 0$. For any infinite $\mathcal{A}\subset \mathbb{N}$ we clearly have
$
d(\mathcal{A},c)\le |c|.
$
We are going to show $e|f$. Assume the contrary, i.e., $e\nmid f$. Take $\mathcal{A}_e=\{|e|m: m\in \mathbb{N}\}\subset \mathbb{N}$ and then
$$
\limsup_{n\rightarrow\infty}d(\mathcal{A}_e,en+f)=0.
$$ 
It follows that
\begin{align*}
\limsup_{n\rightarrow\infty}|d(\mathcal{A}_e,c)-d(\mathcal{A}_e,en+f)|\le |c|<\infty,
\end{align*}
which is a contradiction. ({\bf Case 2.} corresponds with II and III.)

{\bf Case 3.} $be \neq 0$.

In this case, we firstly show that $c^2+f^2\neq 0$, i.e., at least one of $c$ and $f$ is nonzero. Otherwise, we choose a prime $p_0$ so that $p_0\nmid b$ and $p_0\nmid e$ and take 
\begin{align}\label{newadded1}
\mathcal{A}_0=\{p_0^m: m\in \mathbb{N}\}\subset \mathbb{N}.
\end{align} 
Then we have
\begin{align*}
\limsup_{n\rightarrow\infty}|d(\mathcal{A}_0,bn)-d(\mathcal{A}_0,en)|&=\limsup_{n\rightarrow\infty}|d(\mathcal{A}_0,n)-d(\mathcal{A}_0,n)|=0,
\end{align*} 
which is a contradiction. 

Next, we prove $b|c$ or $e|f$. Assume the contrary. Then we may assume $b=xb_1, c=xc_1, e=ye_1$, and $f=yf_1$ with $b_1,e_1\neq 1$, $\gcd(b_1,c_1)=1$, and $\gcd(e_1,f_1)=1$. We take $\mathcal{A}_{b_1,e_1}=\{(b_1e_1)^m: m\in \mathbb{N}\}\subset \mathbb{N}$. We claim that $d(\mathcal{A}_{b_1,e_1},bn+c)$ and $d(\mathcal{A}_{b_1,e_1},en+f)$ are bounded for all $n$, from which we have
$$
\limsup_{n\rightarrow\infty}|d(\mathcal{A}_{b_1,e_1},bn+c)-d(\mathcal{A}_{b_1,e_1},en+f)|<\infty,
$$
leading to a contradiction with (\ref{eq4}). Actually, for all sufficiently large $m$ if $(b_1e_1)^m|bn+c$, then $b_1|b_1n+c_1$. This is a contradiction with $\gcd(b_1,c_1)=1$ and $b_1>1$. Thus, the function $d(\mathcal{A}_{b_1,e_1},bn+c)$ is bounded for all $n$. The same argument yields that $d(\mathcal{A}_{b_1,e_1},en+f)$ is bounded for all $n$.

We are left over to show $bf\neq ec$. Assume the contrary, i.e., 
$$
\frac{c}{b}=\frac{f}{e}=g,\quad \text{say.}
$$
Since $c^2+f^2\neq 0$, we have $g\neq 0$. Moreover, $g\in \mathbb{Z}$ since $b|c$ or $e|f$. Then for $\mathcal{A}_0$ as in (\ref{newadded1}) we also have
\begin{align*}
\limsup_{n\rightarrow\infty}|d(\mathcal{A}_0,bn+c)-d(\mathcal{A}_0,en+f)|&=\limsup_{n\rightarrow\infty}|d(\mathcal{A}_0,b(n+g))-d(\mathcal{A}_0,e(n+g))|\\
&=\limsup_{n\rightarrow\infty}|d(\mathcal{A}_0,n+g)-d(\mathcal{A}_0,n+g)|=0,
\end{align*} 
which is a contradiction. 
({\bf Case 3.} corresponds with IV.)

{\bf Sufficiencies.} We now turn to the proof of sufficiencies case by case. 

{\bf I. implies (\ref{eq4}).} Clearly.

{\bf II. implies (\ref{eq4}).} Without loss of generality, we only consider the case $b=0$, $e\neq 0$, and $c=0$ since the other one can be proved similarly. We have $d(\mathcal{A},en+f)<\infty$ for any $n\neq -f/e$ since $e\neq 0$. But $d(\mathcal{A},0)=\infty$ for any $n$, from which $(\ref{eq4})$ clearly follows.

{\bf III. implies (\ref{eq4}).} Without loss of generality, we only consider the case $b=0$, $e\neq 0$, $c\neq 0$, and $e|f$ since the other one can be proved similarly. We assume $f=eh$. For any positive integer $t$ let
$$
P_t=\prod_{1\le i\le t}a_i.
$$
For any $t$ there is a positive integer $n_t$ so that
$$
n_t\equiv -h\pmod {P_{t}},
$$
which means that 
$
P_t|e(n_t+h).
$
In other words, $d(\mathcal{A},n_t+h)\ge t$.
Hence, for sufficiently large $t$ we have 
\begin{align*}
\limsup_{n\rightarrow\infty}|d(\mathcal{A},c)-d(\mathcal{A}_g,en+f)|&\ge\limsup_{t\rightarrow\infty}|d(\mathcal{A},c)-d(\mathcal{A},e(n_t+h))|\\
&\ge \limsup_{t\rightarrow\infty}(t-|c|)=\infty.
\end{align*} 

{\bf IV. implies (\ref{eq4}).} Without loss of generality, we only consider the case $be \neq 0$, $bf\neq ec$, and $b|c$ since the other one can be proved similarly.

Suppose that $\mathcal{A}=\{a_1<a_2<\cdots<a_i<\cdots\}$ and $c=bh$. Since $bf\neq ec$, we have $f-eh\neq 0$. For any positive integer $t$ let
$$
P_t=\prod_{1\le i\le t}a_i \quad \text{and} \quad \ell=\sup_{t\ge 1}\big\{\gcd(eP_t,f-eh)\big\}.
$$
Then there is some fixed $t_0$ so that $\ell|eP_{t_0}$ and
\begin{align}\label{newadded2}
P_{t_0}>|h|. 
\end{align}
For any positive integer $m$, we  have 
$$
\gcd\left(\frac{eP_{t_0+m}}{\ell},\frac{f-eh}{\ell}\right)=1.
$$
Now, we will separate the arguments into two cases. 

Suppose first that $e>0$.  
In this case, by Dirichlet's theorem in arithmetic progressions there exists a prime 
\begin{align}\label{newadded3}
p_m>\max\Big\{\frac{|f-eh|}{\ell},\ell\Big\}
\end{align}
satisfying that
\begin{align}\label{eq5}
p_m\equiv \frac{f-eh}{\ell} \pmod{eP_{t_0+m}/\ell},
\end{align}
which clearly means that
\begin{align}\label{eq6}
\ell p_m-f+eh=eP_{t_0+m}q_m
\end{align}
for some integer $q_m$. Note that $p_m>\frac{|f-eh|}{\ell}$ from (\ref{newadded3}). We have $\ell p_m-f+eh>0$ and hence $P_{t_0+m}q_m>0$.
Rearranging the terms in (\ref{eq6}) leads to
\begin{align}\label{eq7}
e(P_{t_0+m}q_m-h)+f=\ell p_m.
\end{align}
Moreover, one notes easily that
\begin{align}\label{eq8}
b(P_{t_0+m}q_m-h)+c=bP_{t_0+m}q_m\equiv 0\pmod{P_{t_0+m}}.
\end{align}
From (\ref{eq7}) and (\ref{eq8}) we have
\begin{align}\label{eq9}
d(\mathcal{A},e(P_{t_0+m}q_m-h)+f)=d(\mathcal{A},\ell p_m)
\end{align}
and
\begin{align}\label{eq10}
d(\mathcal{A},b(P_{t_0+m}q_m-h)+c)\ge t_0+m.
\end{align}
Recall that $p_m>\ell$ for any $m$ from (\ref{newadded3}). Hence 
\begin{align}\label{eq11}
d(\mathcal{A},\ell p_m)\le 2\ell\le 2|f-eh|.
\end{align}
Furthermore, from (\ref{newadded2}) we have 
$
P_{t_0+m}q_m-h>P_{t_0}-h>0.
$
We conclude from (\ref{eq9}), (\ref{eq10}) and (\ref{eq11}) that for sufficiently large $m$,
\begin{align*}
\limsup_{n\rightarrow\infty}|&d(\mathcal{A},bn+c)-d(\mathcal{A},en+f)|\\
&\ge\limsup_{m\rightarrow\infty}|d(\mathcal{A},b(P_{t_0+m}q_m-h)+c)-d(\mathcal{A},e(P_{t_0+m}q_m-h)+f)|\\
&=\limsup_{m\rightarrow\infty}|d(\mathcal{A},b(P_{t_0+m}q_m-h)+c)-d(\mathcal{A},\ell p_m)|\\
&\ge \limsup_{m\rightarrow\infty}\big(t_0+m-2|f-eh|\big)=\infty.
\end{align*}

It remains to consider the case $e<0$. In this case, we simply replace (\ref{eq5}) by
$$
p_m\equiv -\frac{f-eh}{\ell} \pmod{eP_{t_0+m}/\ell}.
$$
Then (\ref{eq6}) will be changed to
$$
-\ell p_m-f+eh=eP_{t_0+m}q_m.
$$
We see easily that $q_m$ is also positive. Adjusting the arguments accordingly would yield our desired result. 
\end{proof}

\end{document}
